\documentclass[10pt,a4paper]{amsart}
\usepackage[margin=19mm]{geometry}
\usepackage{amsmath,amssymb,mathtools}
\usepackage[scr=rsfs,cal=euler]{mathalfa}
\usepackage{xfrac}
\usepackage[unicode,hidelinks,bookmarks=false]{hyperref}
\newcommand{\ie}{i.e.\ }
\newcommand{\cf}{cf.\ }
\newcommand{\resp}{respectively\ }
\newcommand{\Subsection}{\subsection}
\providecommand{\nobreakdash}{\nobreak\hbox{-}}
\setlength{\emergencystretch}{2em}
\multlinegap0pt
\usepackage{amssymb}
\usepackage[all]{xy}
\newtheorem{prop}{\ \ \ Proposition}
\title{Sufficient conditions for solvability of linear Diophantine equations, and Frobenius numbers}
\author{Eteri Samsonadze}
\date{Source: arXiv:2408.17266v2 (CC BY 4.0)}
\begin{document}
\maketitle

\noindent\emph{Source and licence.} Sufficient conditions for solvability of linear Diophantine equations, and Frobenius numbers. Version arXiv:2408.17266v2 (CC BY 4.0), distributed by arXiv under the Creative Commons Attribution 4.0 International licence, http://creativecommons.org/licenses/by/4.0/, as recorded in the arXiv licence metadata for this exact version and shown on the abstract page https://arxiv.org/abs/2408.17266v2. Full licence notice: \texttt{licenses/arxiv-2408.17266-CC-BY-4.0.txt}.\par\smallskip

\noindent\textit{Editorial scope.} Counted unit: the article's prop prop (source0.tex lines 381--387). The article's complete original proof of it is delivered with it (source0.tex lines 389--405, one \texttt{proof} environment of 17 lines). No mathematics is written, repaired or re-derived: the statement and the proof are the article's own lines, in the article's own notation, and no retained line is substituted or re-flowed.  No further source-local context is needed: every cross-reference of every retained line resolves inside the two delivered spans.  Nothing was omitted from any retained span, so the package carries no omission record.  No retained line contains a citation, so the wrapper carries no bibliography and no bibliography entry.  No retained line contains a figure environment, an image-inclusion command or drawing code, so no image is delivered and the package declares no asset dependency.  Editorial formatting only, in addition: an AMS \texttt{amsart} preamble that restates the article's own declarations and macros used by the retained lines, namely the environments \texttt{prop} on separate counters, together with the standard packages those lines need.  The retained environments print in this document as Proposition 1 in order of appearance, whereas the article prints the same items with the numbering of its own full document.  The complete native source of the article as distributed in the arXiv e-print for this exact version is bundled unchanged as \texttt{sources/164512/source0.tex}.
\begin{prop}
We have
\begin{equation}
g(a,a+1,a+2,...,g(a,a+1))=a-1
\end{equation}
for $a\geq 3$.
\end{prop}

\begin{proof}
First note that since 
\begin{equation}
g(a,a+1)=a^2-a-1,
\end{equation}
we have
\begin{equation}
g(a,a+1)>a+1
\end{equation}
for $a\geq 3$.

It is obvious that the equation 
\begin{equation}
ax_1+(a+1)x_2+(a+2)x_3+...+g(a,a+1)x_n=b
\end{equation}
is solvable in integer non-negative numbers for $b=a, a+1, a+2,...,g(a,a+1)$. Moreover, since $(a,a+1)=1$, Proposition 2.5 implies that it is solvable for $b>g(a,a+1)$. Thus, this equation is solvable for $b\geq a$. It is obvious that this equation is not solvable for $b=a-1$. This implies (2.16).
\end{proof}

\end{document}
